%!TEX TS-program = xelatex
%!TEX encoding = UTF-8 Unicode
% Awesome CV LaTeX Template for Cover Letter
%
% This template has been downloaded from:
% https://github.com/posquit0/Awesome-CV
%
% Authors:
% Claud D. Park <posquit0.bj@gmail.com>
% Lars Richter <mail@ayeks.de>
%
% Template license:
% CC BY-SA 4.0 (https://creativecommons.org/licenses/by-sa/4.0/)
%


%-------------------------------------------------------------------------------
% CONFIGURATIONS
%-------------------------------------------------------------------------------
% A4 paper size by default, use 'letterpaper' for US letter
\documentclass[11pt, a4paper]{awesome-cv}

% Configure page margins with geometry
\geometry{left=1.8cm, top=.8cm, right=1.8cm, bottom=1.8cm, footskip=.5cm}

% Color for highlights
% Awesome Colors: awesome-emerald, awesome-skyblue, awesome-red, awesome-pink, awesome-orange
%                 awesome-nephritis, awesome-concrete, awesome-darknight
\colorlet{awesome}{awesome-red}
% Uncomment if you would like to specify your own color
% \definecolor{awesome}{HTML}{CA63A8}

% Colors for text
% Uncomment if you would like to specify your own color
% \definecolor{darktext}{HTML}{414141}
\definecolor{text}{HTML}{222222}
\definecolor{graytext}{HTML}{222222}
% \definecolor{lighttext}{HTML}{999999}
% \definecolor{sectiondivider}{HTML}{5D5D5D}

% Set false if you don't want to highlight section with awesome color
\setbool{acvSectionColorHighlight}{false}

% If you would like to change the social information separator from a pipe (|) to something else
\renewcommand{\acvHeaderSocialSep}{\quad\textbar\quad}


%-------------------------------------------------------------------------------
%	PERSONAL INFORMATION
%	Comment any of the lines below if they are not required
%-------------------------------------------------------------------------------
% Available options: circle|rectangle,edge/noedge,left/right
\photo[circle,noedge,right]{../../shared/profile.jpg}
\name{Laurent}{Colbois}
\position{AI / Machine Learning Engineer{\enskip\cdotp\enskip}PhD in Machine Learning \& Computer Vision}
\address{Rue du Rhône 2, 1920 Martigny, Switzerland}

\mobile{+41 79 942 29 61}
\email{laurent.colbois@gmail.com}
\dateofbirth{May 2nd, 1996}
\linkedin{lcolbois}
%\orcid{0009-0006-3547-5841}
%\googlescholar{LOLk6vAAAAAJ}{}

%-------------------------------------------------------------------------------
%	LETTER INFORMATION
%	*** CUSTOMIZE FOR EACH APPLICATION ***
%-------------------------------------------------------------------------------
% The company being applied to
\recipient
  {Hiring Team} % *** CUSTOMIZE: Use name if known, or "Hiring Team" ***
  {BeEmotion\\Lausanne VD, Switzerland} % *** CUSTOMIZE: Company details ***
% The date on the letter, default is the date of compilation
\letterdate{\today}
% The title of the letter
\lettertitle{Application for AI Application Engineer} % *** CUSTOMIZE: Exact job title ***
% How the letter is opened
\letteropening{Dear Hiring Team,} % *** CUSTOMIZE: Use name if known, or "Hiring Team" ***
% How the letter is closed
\letterclosing{Sincerely,}
% Any enclosures with the letter
%\letterenclosure[Attached]{Curriculum Vitae}


%-------------------------------------------------------------------------------
\begin{document}

% Print the header with above personal information
% Give optional argument to change alignment(C: center, L: left, R: right)
\makecvheader[R]

% Print the footer with 3 arguments(<left>, <center>, <right>)
% Leave any of these blank if they are not needed
\makecvfooter
  {\today}
  {Laurent Colbois~~~·~~~Cover Letter}
  {}

% Print the title with above letter information
\makelettertitle

%-------------------------------------------------------------------------------
%	LETTER CONTENT
%-------------------------------------------------------------------------------
\begin{cvletter}
I am applying for the AI Application Engineer position at BeEmotion. After five years building and evaluating machine learning systems at the Idiap Research Institute, I am transitioning towards applied AI engineering, and BeEmotion's focus on taking ML models from prototype to production (training, building inference pipelines, and packaging them for cloud and edge deployment) maps closely onto the direction I want to take. What makes BeEmotion particularly compelling is that your core technology sits directly adjacent to my own research background: my PhD and postdoctoral work centred on deep learning for face recognition, biometrics, and deepfake detection. This experience building and evaluating systems that extract meaningful signals from facial imagery  connects naturally to your work on vision-based behavioral sensing.


My research background at Idiap directly maps onto the technical core of this role. During my postdoc, I built end-to-end evaluation pipelines for Vision-Language Models applied to face recognition tasks, from data preparation through automated evaluation and inference optimization. This built on my PhD work in deep learning and computer vision, where I developed image manipulation detection models and ran systematic, large-scale evaluations of face recognition systems against adversarial attacks, processing large-scale face image datasets. Throughout my work, I apply rigorous MLOps practices: configuration management (Hydra), experiment tracking (MLflow), pipeline orchestration (Snakemake), containerization (Docker/Apptainer), building reproducible workflows which can support the path from prototypes to production-ready systems. Earlier in my career, I also deployed a deep learning model to embedded hardware at Logitech as a working demo, optimizing for real-time inference under latency and memory constraints, which gave me early hands-on exposure to edge deployment constraints. I have not yet worked with AWS, Azure, or GCP in a production context, but my background in Slurm-managed HPC and containerized environments gives me a solid foundation to ramp up quickly, and I am always driven to learn new tools.

I also have a track record of operating effectively under time pressure in small, cross-functional teams. During a 9-day startup hackathon (Idiap Create Challenge 2025), I built a working prototype end-to-end and co-presented the pitch to a technical and business jury, winning first place in a team of three. During my postdoc, I also coordinated a small interdisciplinary project on a forensic data collection platform, where much of my role was taking initially loose, informal requirements and shaping them step by step into a concrete technical specification the team could build on.

In my previous position, a dedicated team handled technology transfer when applicable, while my role centred on research workflows and model benchmarking. I am now looking to be more directly involved in bringing research prototypes to finalized solutions, and BeEmotion's work at the intersection of computer vision and real-world behavioral sensing, deployed at the edge and in the cloud, is exactly the environment in which I would like to grow. I am a native French speaker and fully fluent in English, and I hold Swiss citizenship. My postdoc contract ended in March, so I am available immediately, and I would welcome the opportunity to discuss my application in an interview.


\end{cvletter}


%-------------------------------------------------------------------------------
% Print the signature and enclosures with above letter information
\makeletterclosing

\end{document}
